\chapter{2005年北京卷（春理）}

\section{单选题}

\begin{problem}
\(\i-2\) 的共轭复数是（\quad）

\choices
    {\(2+\i\)}
    {\(2-\i\)}
    {\(-2+\i\)}
    {\(-2-\i\)}
\end{problem}

\begin{answer}
D
\end{answer}

\begin{solution}
复数 \(z=a+b\i\) 的共轭复数为 \(\overline{z}=a-b\i\). 这里 \(\i-2=-2+\i\)，所以其共轭复数为 \(-2-\i\)，故选 D.
\end{solution}

\begin{problem}
函数 \(y=|\log_2 x|\) 的图象是（\quad）

\choices
    {\choicebitmap[width=0.15\paperwidth]{img/2005/beijing_science_spring/q2_optA.png}}
    {\choicebitmap[width=0.15\paperwidth]{img/2005/beijing_science_spring/q2_optB.png}}
    {\choicebitmap[width=0.15\paperwidth]{img/2005/beijing_science_spring/q2_optC.png}}
    {\choicebitmap[width=0.15\paperwidth]{img/2005/beijing_science_spring/q2_optD.png}}
\end{problem}

\begin{answer}
A
\end{answer}

\begin{solution}
函数的定义域为 \(x>0\). 当 \(0<x\leqslant1\) 时，\(\log_2x\leqslant0\)，所以 \(y=-\log_2x\)；当 \(x\geqslant1\) 时，\(y=\log_2x\). 因而图象在直线 \(x=1\) 处取得最小值 \(0\)，且当 \(x\to0^+\) 或 \(x\to+\infty\) 时，\(y\to+\infty\). 与题中图象相符的是 A.
\end{solution}

\begin{problem}
有如下三个命题：

\begin{circlelist}
\item 分别在两个平面内的两条直线一定是异面直线；
\item 垂直于同一个平面的两条直线是平行直线；
\item 过平面 \(\alpha\) 的一条斜线有一个平面与平面 \(\alpha\) 垂直.
\end{circlelist}

其中正确命题的个数为（\quad）

\choices
    {\(0\)}
    {\(1\)}
    {\(2\)}
    {\(3\)}
\end{problem}

\begin{answer}
C
\end{answer}

\begin{solution}
第一个命题不正确：若两个平面 \(\alpha\)，\(\beta\) 相交，取交线上的一点 \(P\)，分别在 \(\alpha\)，\(\beta\) 内作经过 \(P\) 的两条不重合直线，则这两条直线相交，不是异面直线. 第二个命题正确：过空间同一点只有一条直线垂直于已知平面，因此垂直于同一平面的两条直线方向相同，从而平行. 对于第三个命题，设斜线 \(l\) 与平面 \(\alpha\) 交于点 \(P\)，过 \(P\) 作直线 \(m\perp\alpha\)，由 \(l\) 与 \(m\) 确定平面 \(\beta\). 平面 \(\beta\) 含有垂直于平面 \(\alpha\) 的直线 \(m\)，所以 \(\beta\perp\alpha\)，第三个命题正确. 因此正确命题有 \(2\) 个，故选 C.
\end{solution}

\begin{problem}
如果函数 \(f(x)=\sin(\pi x+\theta)\)（\(0<\theta<2\pi\)）的最小正周期是 \(T\)，且当 \(x=2\) 时取得最大值，那么（\quad）

\choices
    {\(T=2\)，\(\theta=\dfrac{\pi}{2}\)}
    {\(T=1\)，\(\theta=\pi\)}
    {\(T=2\)，\(\theta=\pi\)}
    {\(T=1\)，\(\theta=\dfrac{\pi}{2}\)}
\end{problem}

\begin{answer}
A
\end{answer}

\begin{solution}
函数 \(\sin(\pi x+\theta)\) 的最小正周期为 \(T=\dfrac{2\pi}{\pi}=2\). 当 \(x=2\) 取得最大值时，必须有 \(2\pi+\theta=\dfrac\pi2+2k\pi\)，其中 \(k\in\Z\). 因为 \(0<\theta<2\pi\)，只有 \(k=1\) 给出 \(\theta=\dfrac\pi2\). 因此选 A.
\end{solution}

\begin{problem}
设 \(abc\neq0\)，“\(ac>0\)”是“曲线 \(ax^2+by^2=c\) 为椭圆”的（\quad）

\choices
    {充分非必要条件}
    {必要非充分条件}
    {充分必要条件}
    {既非充分又非必要条件}
\end{problem}

\begin{answer}
B
\end{answer}

\begin{solution}
若曲线 \(ax^2+by^2=c\) 是实的非退化椭圆，则 \(a\) 与 \(c\) 同号，且 \(b\) 与 \(c\) 同号，所以必有 \(ac>0\). 但是 \(ac>0\) 不能保证 \(b\) 与 \(c\) 同号，例如 \(a=c=1\)，\(b=-1\) 时曲线为双曲线 \(x^2-y^2=1\). 因此“\(ac>0\)”是必要非充分条件，故选 B.
\end{solution}

\begin{problem}
已知双曲线的两个焦点为 \(F_1\left(-\sqrt5,0\right)\)，\(F_2\left(\sqrt5,0\right)\)，\(P\) 是此双曲线上的一点，且 \(PF_1\perp PF_2\)，\(\left|PF_1\right|\left|PF_2\right|=2\)，则该双曲线的方程是（\quad）

\choices
    {\(\dfrac{x^2}{2}-\dfrac{y^2}{3}=1\)}
    {\(\dfrac{x^2}{3}-\dfrac{y^2}{2}=1\)}
    {\(\dfrac{x^2}{4}-y^2=1\)}
    {\(x^2-\dfrac{y^2}{4}=1\)}
\end{problem}

\begin{answer}
C
\end{answer}

\begin{solution}
设 \(u=\left|PF_1\right|\)，\(v=\left|PF_2\right|\). 因为 \(PF_1\perp PF_2\)，且 \(F_1F_2=2\sqrt5\)，在直角三角形 \(PF_1F_2\) 中有 \(u^2+v^2=20\). 又由题设 \(uv=2\)，所以 \(\left|u-v\right|^2=u^2+v^2-2uv=16\)，从而 \(\left|u-v\right|=4\).

双曲线的焦点在 \(x\) 轴上，设其标准方程为 \(\dfrac{x^2}{a^2}-\dfrac{y^2}{b^2}=1\)，其中 \(a>0\)，焦半距为 \(c=\sqrt5\). 双曲线上任一点到两焦点距离之差的绝对值为 \(2a\)，因此 \(2a=4\)，得 \(a=2\)，\(a^2=4\). 再由 \(c^2=a^2+b^2\) 得 \(b^2=5-4=1\). 所以双曲线方程为 \(\dfrac{x^2}{4}-y^2=1\)，故选 C.
\end{solution}

\begin{problem}
在 \(\triangle ABC\) 中，已知 \(2\sin A\cos B=\sin C\)，那么 \(\triangle ABC\) 一定是（\quad）

\choices
    {直角三角形}
    {等腰三角形}
    {等腰直角三角形}
    {正三角形}
\end{problem}

\begin{answer}
B
\end{answer}

\begin{solution}
在三角形中 \(A+B+C=\pi\)，所以 \(\sin C=\sin(A+B)\). 由题设，\(2\sin A\cos B=\sin(A+B)\)，即 \(2\sin A\cos B=\sin A\cos B+\cos A\sin B\)，从而 \(\sin(A-B)=0\). 由于 \(A\)，\(B\in(0,\pi)\)，有 \(A-B\in(-\pi,\pi)\)，故只能是 \(A=B\). 因而三角形 \(ABC\) 一定是等腰三角形，故选 B.
\end{solution}

\begin{problem}
若不等式 \((-1)^na<2+\dfrac{(-1)^{n+1}}{n}\) 对于任意正整数 \(n\) 恒成立，则实数 \(a\) 的取值范围是（\quad）

\choices
    {\(\left[-2,\dfrac32\right)\)}
    {\(\left(-2,\dfrac32\right)\)}
    {\(\left[-3,\dfrac32\right)\)}
    {\(\left(-3,\dfrac32\right)\)}
\end{problem}

\begin{answer}
A
\end{answer}

\begin{solution}
当 \(n\) 为偶数时，原不等式化为 \(a<2-\dfrac1n\). 所有正偶数中最小的是 \(n=2\)，且当 \(n\geqslant2\) 时有 \(2-\dfrac1n\geqslant\dfrac32\)，因而偶数情形等价于 \(a<\dfrac32\).

当 \(n\) 为奇数时，原不等式化为 \(-a<2+\dfrac1n\)，即 \(a>-2-\dfrac1n\). 若 \(a\geqslant-2\)，则对任意正奇数 \(n\)，都有 \(-2-\dfrac1n<-2\leqslant a\)，所以奇数情形成立. 反之，若 \(a<-2\)，取正奇数 \(n>\dfrac1{-2-a}\)，则 \(\dfrac1n<-2-a\)，从而 \(-2-\dfrac1n>a\)，奇数情形不成立. 因此奇数情形等价于 \(a\geqslant-2\).

综合两种情形，得 \(a\in\left[-2,\dfrac32\right)\).
当 \(a=-2\) 时奇数情形中的不等式仍为严格成立，当 \(a=\dfrac32\) 时 \(n=2\) 不能成立，故端点取舍正确，对应选项为 A.
\end{solution}

\section{填空题}

\begin{problem}
\(\displaystyle\lim_{n\to\infty}\frac{n^2+2n}{2n^2-3}=\fillinblank{}\).
\end{problem}

\begin{answer}
\(\dfrac12\)
\end{answer}

\begin{solution}
将分子、分母同除以 \(n^2\)，得 \(\dfrac{n^2+2n}{2n^2-3}=\dfrac{1+2/n}{2-3/n^2}\). 当 \(n\to\infty\) 时，\(2/n\to0\)，\(3/n^2\to0\)，且分母的极限为 \(2\neq0\)，所以原极限为 \(\dfrac12\).
\end{solution}

\begin{problem}
已知 \(\sin\dfrac{\theta}{2}+\cos\dfrac{\theta}{2}=\dfrac{2\sqrt3}{3}\)，那么 \(\sin\theta\) 的值为\fillinblank{}，\(\cos 2\theta\) 的值为\fillinblank{}.
\end{problem}

\begin{answer}
\(\dfrac13\)，\(\dfrac79\)
\end{answer}

\begin{solution}
令 \(u=\sin\dfrac\theta2\)，\(v=\cos\dfrac\theta2\). 由 \(u+v=\dfrac{2\sqrt3}{3}\)，平方得 \(u^2+v^2+2uv=\dfrac43\). 因为 \(u^2+v^2=1\)，且 \(2uv=\sin\theta\)，所以 \(1+\sin\theta=\dfrac43\)，即 \(\sin\theta=\dfrac13\). 于是 \(\cos2\theta=1-2\sin^2\theta=1-\dfrac29=\dfrac79\).
\end{solution}

\begin{problem}
若圆 \(x^2+y^2+mx-\dfrac14=0\) 与直线 \(y=-1\) 相切，且其圆心在 \(y\) 轴的左侧，则 \(m=\fillinblank{}\).
\end{problem}

\begin{answer}
\(\sqrt3\)
\end{answer}

\begin{solution}
将圆的方程配方，得 \(\left(x+\dfrac m2\right)^2+y^2=\dfrac{m^2+1}{4}\). 因此圆心为 \(\left(-\dfrac m2,0\right)\)，半径为 \(\dfrac{\sqrt{m^2+1}}2\). 直线 \(y=-1\) 与圆相切，故圆心到该直线的距离等于半径，即 \(1=\dfrac{\sqrt{m^2+1}}2\)，从而 \(m^2=3\)，得 \(m=\pm\sqrt3\). 又圆心在 \(y\) 轴左侧，故 \(-\dfrac m2<0\)，所以 \(m=\sqrt3\).
\end{solution}

\begin{problem}
如图，正方体 \(ABCD-A_1B_1C_1D_1\) 的棱长为 \(a\)，将该正方体沿对角面 \(BB_1D_1D\) 切成两块，再将这两块拼接成一个不是正方体的四棱柱，那么所得四棱柱的全面积为\(\fillinblank{}\).

\bitmapfigure[max width=0.35\linewidth]{img/2005/beijing_science_spring/q12_fig1.png}
\end{problem}

\begin{answer}
\(\left(4+2\sqrt2\right)a^2\)
\end{answer}

\begin{solution}
沿矩形对角面 \(BB_1D_1D\) 切开后，正方体分成两个全等的直三棱柱. 每个三棱柱的两个三角形底面都是直角等腰三角形，面积各为 \(\dfrac{a^2}{2}\)；其三个侧面中有两个边长为 \(a\) 的正方形，另一个是边长为 \(a\) 和 \(\sqrt2a\) 的矩形. 因此每个三棱柱的表面积为 \(a^2+2a^2+\sqrt2a^2=(3+\sqrt2)a^2\)，两块的表面积总和为 \((6+2\sqrt2)a^2\).

拼接成四棱柱时，可将两块的一个边长为 \(a\) 的正方形侧面完全重合，该接触面面积为 \(a^2\)，它不属于拼接后四棱柱的外表面，且在两块表面积总和中被计算了两次. 所以所得四棱柱的全面积为 \((6+2\sqrt2)a^2-2a^2=(4+2\sqrt2)a^2\).
\end{solution}

\begin{problem}
从 \(-1\)，\(0\)，\(1\)，\(2\) 这四个数中选三个不同的数作为函数 \(f(x)=ax^2+bx+c\) 的系数，可组成不同的二次函数共有\fillinblank{}个，其中不同的偶函数共有\fillinblank{}个. （用数字作答）
\end{problem}

\begin{answer}
\(18\)，\(6\)
\end{answer}

\begin{solution}
因为函数是二次函数，所以二次项系数 \(a\neq0\)，可从 \(-1\)，\(1\)，\(2\) 中选取，共有 \(3\) 种选法. 选定 \(a\) 后，从剩余三个数中依次选取互不相同的 \(b\)，\(c\)，有 \(3\times2\) 种选法，因此不同的二次函数共有 \(3\times3\times2=18\) 个.

函数为偶函数的充要条件是 \(b=0\). 此时 \(a\) 和 \(c\) 需从 \(-1\)，\(1\)，\(2\) 中选取两个不同的数并分别放在两个位置上，共有 \(3\times2=6\) 个.
\end{solution}

\begin{problem}
若关于 \(x\) 的不等式 \(x^2-ax-a>0\) 的解集为 \(\left(-\infty,+\infty\right)\)，则实数 \(a\) 的取值范围是\fillinblank{}；若关于 \(x\) 的不等式 \(x^2-ax-a\leqslant-3\) 的解集不是空集，则实数 \(a\) 的取值范围是\fillinblank{}.
\end{problem}

\begin{answer}
\(\left(-4,0\right)\)，\(\left(-\infty,-6\right]\cup\left[2,+\infty\right)\)
\end{answer}

\begin{solution}
对于不等式 \(x^2-ax-a>0\)，其二次项系数为正. 要使解集为全体实数，二次函数的最小值必须严格大于 \(0\)，等价于判别式小于 \(0\). 因而 \(\Delta=a^2+4a=a(a+4)<0\)，解得 \(-4<a<0\).

对于不等式 \(x^2-ax-a\leqslant-3\)，移项得 \(x^2-ax+3-a\leqslant0\). 这是开口向上的二次函数，不等式有非空解集的充要条件是判别式不小于 \(0\)，所以 \(\Delta=a^2-4(3-a)=a^2+4a-12=(a+6)(a-2)\geqslant0\). 解得 \(a\leqslant-6\) 或 \(a\geqslant2\)，即 \(a\in\left(-\infty,-6\right]\cup\left[2,+\infty\right)\).
\end{solution}

\section{解答题}

\begin{problem}
设函数 \(f(x)=\lg(2x-3)\) 的定义域为集合 \(M\)，函数 \(g(x)=\sqrt{1-\dfrac{2}{x-1}}\) 的定义域为集合 \(N\). 求：

\begin{enumerate}
\item 集合 \(M\)，\(N\)；
\item 集合 \(M\cap N\)，\(M\cup N\).
\end{enumerate}
\end{problem}

\begin{solution}
\begin{enumerate}
\item 由对数函数的真数必须为正，得 \(2x-3>0\)，所以
\[
M=\left(\dfrac32,+\infty\right)
\]
对于函数 \(g\)，根式内的式子必须非负，且分式的分母不能为零，即 \(x\ne1\) 且 \(1-\dfrac{2}{x-1}\geqslant0\). 将不等式化为
\[
\dfrac{x-3}{x-1}\geqslant0
\]
在关键点 \(1\)，\(3\) 将实数轴分成三个区间. 当 \(x<1\) 时，分子、分母均为负，分式为正；当 \(1<x<3\) 时，分子为负、分母为正，分式为负；当 \(x>3\) 时，分子、分母均为正，分式为正. 又因为 \(x=1\) 不在定义域内、\(x=3\) 使分式为零且可以取到，所以
\[
N=\left(-\infty,1\right)\cup\left[3,+\infty\right)
\]

\item 由 \(M=\left(\dfrac32,+\infty\right)\) 与 \(N=\left(-\infty,1\right)\cup\left[3,+\infty\right)\)，可得 \(M\cap N=\left[3,+\infty\right)\)，\(M\cup N=\left(-\infty,1\right)\cup\left(\dfrac32,+\infty\right)\).
\end{enumerate}
\end{solution}

\begin{problem}
如图，正三角形 \(ABC\) 的边长为 \(3\)，过其中心 \(G\) 作 \(BC\) 边的平行线，分别交 \(AB\)，\(AC\) 于 \(B_1\)，\(C_1\). 将 \(\triangle AB_1C_1\) 沿 \(B_1C_1\) 折起到 \(\triangle A_1B_1C_1\) 的位置，使点 \(A_1\) 在平面 \(BB_1C_1C\) 上的射影恰是线段 \(BC\) 的中点 \(M\). 求：

\begin{enumerate}
\item 二面角 \(A_1-B_1C_1-M\) 的大小；
\item 异面直线 \(A_1B_1\) 与 \(CC_1\) 所成角的大小（用反三角函数表示）.
\end{enumerate}

\bitmapfigure[max width=0.42\linewidth]{img/2005/beijing_science_spring/q16_fig1.png}
\end{problem}

\begin{solution}
以 \(M\) 为原点，取 \(BC\) 所在直线为 \(x\) 轴，\(MG\) 所在直线为 \(y\) 轴，并取垂直于平面 \(ABC\) 的方向为 \(z\) 轴. 由于 \(ABC\) 是边长为 \(3\) 的正三角形，可取 \(B\left(-\dfrac32,0,0\right)\)，\(C\left(\dfrac32,0,0\right)\)，\(A\left(0,\dfrac{3\sqrt3}{2},0\right)\).
正三角形的中心在中线 \(AM\) 上，且 \(AG:GM=2:1\)，所以 \(GM=\dfrac{\sqrt3}{2}\). 直线 \(B_1C_1\) 过 \(G\) 且平行于 \(BC\)，由相似三角形得 \(B_1C_1=2\)，从而 \(AB_1=AC_1=2\). 折叠不改变三角形的边长，故 \(A_1B_1=A_1C_1=2\).

由 \(B_1C_1\parallel BC\) 且 \(B_1C_1=2\)，可得 \(B_1\left(-1,\dfrac{\sqrt3}{2},0\right)\)，\(C_1\left(1,\dfrac{\sqrt3}{2},0\right)\). 又因为点 \(A_1\) 在平面 \(BB_1C_1C\) 上的射影为 \(M\)，所以 \(A_1M\perp\) 平面 \(BB_1C_1C\). 在直角三角形 \(A_1MB_1\) 中，\(MB_1=\dfrac{\sqrt7}{2}\)，于是 \(A_1M=\dfrac32\). 因此 \(A_1G^2=A_1M^2+GM^2=3\)，即 \(A_1G=\sqrt3=2GM\).

\begin{enumerate}
\item 过点 \(G\) 作平面 \(A_1GM\)，它垂直于棱 \(B_1C_1\)，并且与两个面分别交于 \(GA_1\) 和 \(GM\). 因而二面角 \(A_1-B_1C_1-M\) 的平面角为 \(\angle A_1GM\). 在直角三角形 \(A_1GM\) 中，
\[
\cos\angle A_1GM=\dfrac{GM}{A_1G}=\dfrac12
\]
所以二面角 \(A_1-B_1C_1-M\) 的大小为 \(60^\circ\).

\item 过 \(B_1\) 作 \(B_1P\parallel C_1C\)，交 \(BC\) 于 \(P\). 则异面直线 \(A_1B_1\) 与 \(CC_1\) 所成的角等于 \(\angle A_1B_1P\). 由前面的坐标可知 \(B_1\left(-1,\dfrac{\sqrt3}{2},0\right)\)，\(C_1\left(1,\dfrac{\sqrt3}{2},0\right)\)，\(C\left(\dfrac32,0,0\right)\)，所以 \(P\left(-\dfrac12,0,0\right)\). 因而
\[
A_1P^2=A_1M^2+MP^2=\left(\dfrac32\right)^2+\left(\dfrac12\right)^2=\dfrac52
\]
且 \(B_1P=C_1C=1\). 在三角形 \(A_1B_1P\) 中，利用余弦定理得
\[
\cos\angle A_1B_1P
=\dfrac{A_1B_1^2+B_1P^2-A_1P^2}{2A_1B_1\cdot B_1P}
=\dfrac{4+1-\dfrac52}{2\cdot2\cdot1}
=\dfrac58
\]
所以异面直线 \(A_1B_1\) 与 \(CC_1\) 所成角的大小为 \(\arccos\dfrac58\).
\end{enumerate}
\end{solution}

\begin{problem}
已知 \(\{a_n\}\) 是等比数列，\(a_1=2\)，\(a_3=18\)；\(\{b_n\}\) 是等差数列，\(b_1=2\)，\(b_1+b_2+b_3+b_4=a_1+a_2+a_3>20\).

\begin{enumerate}
\item 求数列 \(\{b_n\}\) 的通项公式；
\item 求数列 \(\{b_n\}\) 的前 \(n\) 项和 \(S_n\) 的公式；
\item 设 \(P_n=b_1+b_4+b_7+\cdots+b_{3n-2}\)，\(Q_n=b_{10}+b_{12}+b_{14}+\cdots+b_{2n+8}\)，其中 \(n=1\)，\(2\)，\(\cdots\)，试比较 \(P_n\) 与 \(Q_n\) 的大小，并证明你的结论.
\end{enumerate}
\end{problem}

\begin{solution}
设等比数列 \(\{a_n\}\) 的公比为 \(q\). 由 \(a_3=a_1q^2\)，得 \(2q^2=18\)，即 \(q=\pm3\). 又因为
\[
a_1+a_2+a_3=2+2q+18=20+2q>20
\]
所以 \(q=3\)，从而 \(a_1+a_2+a_3=26\). 设等差数列 \(\{b_n\}\) 的公差为 \(d\)，由题设等式得
\[
b_1+b_2+b_3+b_4=2+(2+d)+(2+2d)+(2+3d)=8+6d=26
\]
故 \(d=3\).

\begin{enumerate}
\item 数列 \(\{b_n\}\) 的通项公式为
\[
b_n=b_1+(n-1)d=2+3(n-1)=3n-1
\]

\item 前 \(n\) 项和为
\[
S_n=\dfrac{n(b_1+b_n)}{2}=\dfrac{n(2+3n-1)}{2}=\dfrac{n(3n+1)}{2}
\]

\item \(P_n\) 的各项下标为 \(1\)，\(4\)，\(7\)，\(\ldots\)，\(3n-2\)，对应的数列是首项 \(b_1=2\)、公差 \(3d\) 的等差数列，所以
\[
P_n=\sum_{k=1}^n b_{3k-2}=\sum_{k=1}^n(9k-7)=\dfrac{n(9n-5)}{2}
\]
\(Q_n\) 的各项下标为 \(10\)，\(12\)，\(14\)，\(\ldots\)，\(2n+8\)，对应的数列是首项 \(b_{10}=2+9d\)、公差 \(2d\) 的等差数列，所以
\[
Q_n=\sum_{k=1}^n b_{2k+8}=\sum_{k=1}^n(6k+23)=n(3n+26)
\]
于是
\[
P_n-Q_n=\dfrac{3n(n-19)}{2}
\]
由于 \(n\) 为正整数，因此当 \(1\leqslant n\leqslant18\) 时，\(P_n<Q_n\)；当 \(n=19\) 时，\(P_n=Q_n\)；当 \(n\geqslant20\) 时，\(P_n>Q_n\).
\end{enumerate}
\end{solution}

\begin{problem}
如图，\(O\) 为坐标原点，直线 \(l\) 在 \(x\) 轴和 \(y\) 轴上的截距分别是 \(a\) 和 \(b\)，且交抛物线 \(y^2=2px\)（\(p>0\)）于 \(M\left(x_1,y_1\right)\)，\(N\left(x_2,y_2\right)\) 两点.

\begin{enumerate}
\item 写出直线 \(l\) 的截距式方程；
\item 证明：\(\dfrac{1}{y_1}+\dfrac{1}{y_2}=\dfrac{1}{b}\)；
\item 当 \(a=2p\) 时，求 \(\angle MON\) 的大小.
\end{enumerate}

\bitmapfigure[max width=0.42\linewidth]{img/2005/beijing_science_spring/q18_fig1.png}
\end{problem}

\begin{solution}
\begin{enumerate}
\item 直线 \(l\) 在坐标轴上的截距分别为 \(a\)，\(b\)，所以其截距式方程为
\[
\dfrac{x}{a}+\dfrac{y}{b}=1
\]

\item 由题设直线有两个坐标轴截距，且与抛物线交于两点，故 \(a\ne0\)，\(b\ne0\). 由上式得 \(x=a\left(1-\dfrac{y}{b}\right)\). 代入抛物线方程 \(y^2=2px\)，得到
\[
y^2+\dfrac{2pa}{b}y-2pa=0
\]
由于 \(M\)，\(N\) 在抛物线上，且它们的纵坐标为该方程的两个根 \(y_1\)，\(y_2\)，由韦达定理得
\(y_1+y_2=-\dfrac{2pa}{b}\)，
\(y_1y_2=-2pa\).
注意到 \(p>0\)，且截距 \(a\ne0\)，所以 \(y_1y_2\ne0\). 因而
\[
\dfrac1{y_1}+\dfrac1{y_2}=\dfrac{y_1+y_2}{y_1y_2}=\dfrac1b
\]

\item 当 \(a=2p\) 时，由上面的根的关系有 \(y_1y_2=-4p^2\). 抛物线上任一点 \((x_i,y_i)\) 满足 \(x_i=\dfrac{y_i^2}{2p}\)，所以直线 \(OM\)，\(ON\) 的斜率分别为
\(k_1=\dfrac{y_1}{x_1}=\dfrac{2p}{y_1}\)，
\(k_2=\dfrac{y_2}{x_2}=\dfrac{2p}{y_2}\).
于是 \(k_1k_2=\dfrac{4p^2}{y_1y_2}=-1\)，故 \(OM\perp ON\). 因而 \(\angle MON=90^\circ\).
\end{enumerate}
\end{solution}

\begin{problem}
经过长期观测得到：在交通繁忙的时段内，某公路段汽车的车流量 \(y\)（\(\text{千辆/小时}\)）与汽车的平均速度 \(v\)（\(\text{千米/小时}\)）之间的函数关系为：
\(y=\frac{920v}{v^2+3v+1600}\)（\(v>0\)）.

\begin{enumerate}
\item 在该时段内，当汽车的平均速度 \(v\) 为多少时，车流量最大？最大车流量为多少？（精确到 \(0.1\text{ 千辆/小时}\)）
\item 若要求在该时段内车流量超过 \(10\text{ 千辆/小时}\)，则汽车的平均速度应在什么范围内？
\end{enumerate}
\end{problem}

\begin{solution}
\begin{enumerate}
\item 由题意，\(v>0\)，且
\[
y(v)=\dfrac{920v}{v^2+3v+1600}
\]
由于 \(v^2+3v+1600>0\)，求导得
\[
y'(v)=\dfrac{920(1600-v^2)}{(v^2+3v+1600)^2}
\]
因此当 \(0<v<40\) 时，\(y'(v)>0\)；当 \(v>40\) 时，\(y'(v)<0\). 函数在 \((0,40)\) 上递增、在 \((40,+\infty)\) 上递减，所以在定义域 \(v>0\) 内于 \(v=40\text{ 千米/小时}\) 处取得全局最大值，且
\[
y(40)=\dfrac{920\times40}{40^2+3\times40+1600}=\dfrac{920}{83}\approx11.1
\]
故最大车流量约为 \(11.1\text{ 千辆/小时}\).

\item 车流量超过 \(10\text{ 千辆/小时}\) 时，由分母为正，得
\[
\dfrac{920v}{v^2+3v+1600}>10
\Longleftrightarrow v^2-89v+1600<0
\]
因式分解为 \((v-25)(v-64)<0\)，再结合 \(v>0\)，可得汽车的平均速度范围为
\[
25<v<64
\]
其中 \(v\) 的单位为 \(\text{千米/小时}\).
\end{enumerate}
\end{solution}

\begin{problem}
现有一组互不相同且从小到大排列的数据：\(a_0\)，\(a_1\)，\(a_2\)，\(a_3\)，\(a_4\)，\(a_5\)，其中 \(a_0=0\). 为提取反映数据间差异程度的某种指标，今对其进行如下加工：记 \(T=a_0+a_1+\cdots+a_5\)，\(x_n=\frac n5\)，\(y_n=\frac1T\left(a_0+a_1+\cdots+a_n\right)\)，作函数 \(y=f(x)\)，使其图象为逐点依次连接点 \(P_n\left(x_n,y_n\right)\)（\(n=0\)，\(1\)，\(2\)，\(\ldots\)，\(5\)）的折线.

\begin{enumerate}
\item 求 \(f(0)\) 和 \(f(1)\) 的值；
\item 设 \(P_{n-1}P_n\) 的斜率为 \(k_n\)（\(n=1\)，\(2\)，\(3\)，\(4\)，\(5\)），判断 \(k_1\)，\(k_2\)，\(k_3\)，\(k_4\)，\(k_5\) 的大小关系；
\item 证明：当 \(x\in\left(0,1\right)\) 时，\(f(x)<x\)；
\item 求由函数 \(y=x\) 与 \(y=f(x)\) 的图象所围成图形的面积（用 \(a_1\)，\(a_2\)，\(a_3\)，\(a_4\)，\(a_5\) 表示）.
\end{enumerate}
\end{problem}

\begin{solution}
由 \(a_0=0<a_1<a_2<a_3<a_4<a_5\)，可知 \(T=a_0+a_1+\cdots+a_5>0\). 记 \(A_n=a_1+a_2+\cdots+a_n\)（\(n=1\)，\(2\)，\(\ldots\)，\(5\)），则 \(T=A_5\)，并且 \(y_n=\dfrac{A_n}{T}\)，其中 \(y_0=0\).

\begin{enumerate}
\item 由 \(P_0=(0,0)\)，\(P_5=(1,1)\)，可知
\(f(0)=0\)，\(f(1)=1\).

\item 由相邻两点的横坐标之差为 \(\dfrac15\)，得
\[
k_n=\dfrac{y_n-y_{n-1}}{x_n-x_{n-1}}=\dfrac{a_n/T}{1/5}=\dfrac{5a_n}{T}
\]
因 \(a_1<a_2<a_3<a_4<a_5\)，所以
\[
k_1<k_2<k_3<k_4<k_5
\]

\item 对任意 \(n=1\)，\(2\)，\(3\)，\(4\)，当 \(1\leqslant i\leqslant n<j\leqslant5\) 时都有 \(a_i<a_j\). 将这 \(n(5-n)\) 个严格不等式逐项相加，得到
\[
\sum_{i=1}^n\sum_{j=n+1}^5 a_i<\sum_{i=1}^n\sum_{j=n+1}^5 a_j
\]
从而
\[
(5-n)A_n<n(T-A_n)
\]
由于 \(n(5-n)>0\)，这等价于
\[
\dfrac{A_n}{n}<\dfrac{T-A_n}{5-n}
\]
整理得 \(5A_n<nT\)，从而
\[
y_n=\dfrac{A_n}{T}<\dfrac n5=x_n
\]
因此折线的各个内点 \(P_1\)，\(P_2\)，\(P_3\)，\(P_4\) 均在直线 \(y=x\) 的下方，而 \(P_0\)，\(P_5\) 在直线 \(y=x\) 上. 在每个区间 \(\left[x_{n-1},x_n\right]\) 内，\(f(x)\) 与 \(x-f(x)\) 都是线性函数，端点处差值非负且不可能在整个区间恒为零，所以当 \(x\in\left(0,1\right)\) 时，\(f(x)<x\).

\item 因为在 \(\left(0,1\right)\) 内有 \(f(x)<x\)，所求面积为
\[
S=\int_0^1\left(x-f(x)\right)\,\mathrm{d}x=\dfrac12-\int_0^1f(x)\,\mathrm{d}x
\]
折线函数的积分等于各小区间上梯形面积之和，所以
\[
\int_0^1f(x)\,\mathrm{d}x=\dfrac15\left(\dfrac{y_0+y_5}{2}+y_1+y_2+y_3+y_4\right)=\dfrac1{10}+\dfrac{4a_1+3a_2+2a_3+a_4}{5T}
\]
代入 \(T=a_1+a_2+a_3+a_4+a_5\)，得
\[
S=\dfrac{2a_5+a_4-a_2-2a_1}{5\left(a_1+a_2+a_3+a_4+a_5\right)}
\]
\end{enumerate}
\end{solution}
